\chapter{A Short History}
\label{app:history}

The methods in this book were invented four times, in four fields, over fifty years,
largely without cross-reference. This appendix tells that story, because knowing it
changes how one reads the literature: a graduate student who meets least angle
regression as a statistical innovation and the Critical Line Algorithm as a piece of
finance folklore will not think to look for the third and fourth copies.

The account is a sketch, not a survey. Each of the four literatures has a history deep
enough to fill an article of its own.

\section{The 1950s: optimization and finance}

The active-set treatment of the parametric quadratic program is a creation of the
1950s, and the Critical Line Algorithm is its earliest fully worked instance.

The decade was the founding one for operations research, and its founding algorithm was
Dantzig's simplex method \cite{dantzig1951}, whose vertex-to-vertex walk along the edges
of a polyhedron the homotopy of Chapter~\ref{ch:parametric} still rides---the ratio test
and the Bland tie-break of that chapter are literally the simplex method's. The same
decade supplied the optimality conditions on which everything in this book rests:
Kuhn and Tucker \cite{kuhntucker1951}, anticipated by Karush \cite{karush1939},
characterised the constrained optimum by the stationarity and complementary-slackness
relations whose parametric solution is the bordered system of
Lemma~\ref{lem:affine}.

Markowitz worked alongside Dantzig at the RAND Corporation and took the method as his
starting point. Having posed mean--variance selection \cite{markowitz1952}, he
recognised that the efficient set is obtained by solving a quadratic program
\emph{parametrically} in the risk--return trade-off, and \cite{markowitz1956} gave the
Critical Line Algorithm, developed in full in the 1959 monograph \cite{markowitz1959}.

The idea did not appear in a vacuum. Gass and Saaty \cite{gasssaaty1955} introduced
\emph{parametric} linear programming, the direct linear ancestor: follow the optimal
vertices as a cost coefficient varies, make the objective quadratic, and that trajectory
becomes the critical line. Frank and Wolfe \cite{frankwolfe1956}, Beale
\cite{beale1959}, and Wolfe \cite{wolfe1959} gave general quadratic-programming
algorithms over the same years, several appearing in the same journal as Markowitz's---
the Critical Line Algorithm contemporary with, and independent of, the general theory
that would later be seen to contain it.

Two threads then ran from this beginning for half a century, both inside finance. On the
modelling side, Sharpe \cite{sharpe1963}, again at Markowitz's suggestion, introduced the
single-index model, writing the covariance as a diagonal matrix plus a rank-one term---
precisely the structure Chapter~\ref{ch:operator} exploits, and the ancestor of every
modern multi-factor risk model. On the algorithmic side the CLA was re-derived, refined,
and scaled across decades, almost always from within portfolio theory
\cite{perold1984,hirschberger2010,niedermayer2010,markowitz2000,stein2008,bailey2013,markowitz2021}.
The lineage is unbroken; what it never asked was whether the algorithm belonged to a
family larger than portfolio selection.

\begin{remark}[Why the method never travelled]
One can only speculate, but the honest reading is that it was easy to do without. A
single efficient point is a quadratic program a generic solver dispatches in
milliseconds, and one approximates the frontier by re-solving as the return tilt is
varied. The parametric machinery that returns the whole frontier exactly bought little
that a loop over a black-box solver did not.

The surviving implementations bear this out: the canonical reference codes support
neither a factor risk model in place of a dense covariance nor inequality constraints
beyond simple bounds---precisely the generality Chapter~\ref{ch:operator} draws out.

The neglect cuts close to home. Markowitz's 1952 paper is among the most-cited articles
in all of economics, while the 1956 paper that actually gave the algorithm is
comparatively forgotten, with no clean citation record of its own, references scattering
across the journal article, its 1955 RAND memorandum, and later reprints. If finance
largely overlooked its own computational contribution, it can hardly fault statistics for
later arriving at the same homotopy without reference to it.
\end{remark}

\section{The 1990s and 2000s: statistics}

Statistics arrived at the same device around forty years later, by a different door, and
at first without reference to any of this.

The door was variable selection and shrinkage, not optimization. The question was not how
to trace a frontier but how the LASSO of Tibshirani \cite{tibshirani1996} chooses, as its
penalty relaxes, which coefficients to admit---and the discovery was that it does so
along a piecewise-linear path. Osborne, Presnell and Turlach \cite{osborne2000} gave the
homotopy explicitly; Efron, Hastie, Johnstone and Tibshirani \cite{efron2004} gave least
angle regression, read the path's corners as the events of a single forward procedure,
and made the piecewise-linear path famous well beyond the LASSO itself. Hastie et al.\
\cite{hastie2004svm} carried the construction to the support vector machine; Tibshirani
and Taylor \cite{tibshirani2011} extended it to the generalized lasso, with the fused
lasso and trend filtering among its special cases, by a dual path of the same kind.

What had been an optimization fact about parametric quadratic programs was being rebuilt,
theorem by theorem, as a native part of regularised estimation.

The statistical development reached even the governing principle. Rosset and Zhu
\cite{rosset2007} characterised, within statistics, exactly when a regularised solution
path is piecewise linear---a quadratic loss against a polyhedral penalty or
constraint---which is the criterion the parametric-programming literature had carried
since the 1950s, now restated in the field's own vocabulary. Tibshirani
\cite{tibshirani2013} settled when the LASSO solution, and so the path, is unique.

None of this leaned on the optimization lineage, and there was little reason it should
have: the statistical versions answered statistical questions---about selection, degrees
of freedom, and prediction---that the portfolio and control literatures had never posed.
That the underlying object was the same went unremarked because no one had occasion to
remark it.

The bridge back was late and partial. Brodie et al.\ \cite{brodie2009} observed that a
gross-exposure-constrained Markowitz problem \emph{is} a LASSO---the special case of
Theorem~\ref{thm:identity} without general constraints---and Gaines, Kim and Zhou
\cite{gaines2018} later traced the constrained LASSO path directly, under arbitrary
linear equalities and inequalities, by an active-set homotopy on the same KKT system.
Even so, the observation travelled as a curiosity rather than as the statement that two
large algorithmic literatures describe one structure.

\section{The 2000s: control, and a partial crossing}

The same device was found once more, in control, and this time the crossing was partly
deliberate.

Through the 2000s the constrained linear-quadratic regulator was being solved
\emph{explicitly}, as a function of the state rather than afresh at each sampling
instant. Bemporad et al.\ \cite{bemporad2002} showed that the optimal control law is
piecewise affine over a polyhedral partition of the state space, and Tøndel, Johansen and
Bemporad \cite{tondel2003} gave the multi-parametric quadratic-programming algorithm
computing it by traversing those regions. This is the common scheme carrying a
\emph{vector} parameter (Remark~\ref{rem:vector}).

Roll \cite{roll2008} then made the link to the scalar case plain, tracing
single-parameter piecewise-linear solution paths in the control setting and explicitly
generalising the LARS/LASSO construction. Control thus arrived at both the scalar
homotopy and its multi-parameter form, and arrived---in Roll's case---already aware of
where statistics had been: the lone instance in this story of a field reaching the device
knowing the others had been there.

\section{The single-point solvers}

The parametric story is the striking one, but the single-point methods of Part~II have
their own chronology and it is worth recording.

The dual active-set method is Goldfarb and Idnani
\cite{goldfarb1982,goldfarb1983}. Its restriction to positive definite Hessians was
lifted by Boland \cite{boland1997}; Forsgren, Gill and Wong \cite{forsgren2016} proved
finite termination for paired primal and dual active-set methods more generally; and
Arnström, Bemporad and Axehill \cite{arnstrom2022daqp} rebuilt the same dual walk on
recursive $LDL\T$ updates for embedded model predictive control---the point at which the
control and optimization threads of this appendix rejoin. Arnström and Axehill
\cite{arnstrom2022certification} additionally compute exact worst-case iteration counts
for parametric QP families.

The complementarity thread is older and separate. Non-negative least squares is Lawson
and Hanson \cite{lawson1974}; principal pivoting and the least-index rule are Murty
\cite{murty1988}; block principal pivoting with the anti-cycling guard of
Theorem~\ref{thm:termination} is Júdice and Pires \cite{judice1994}, with
\cite{cottle1992} the standard reference for linear complementarity. The primal--dual
active set strategy read as semismooth Newton is Hintermüller, Ito and Kunisch
\cite{hintermuller2002}. The Krylov side is Hestenes and Stiefel \cite{hestenes1952},
and the feasible-point family that Remark~\ref{rem:feasible} contrasts runs from
Bertsekas \cite{bertsekas1982} through Moré and Toraldo \cite{moretoraldo1991} to
Dostál \cite{dostal1997,dostalschoberl2005}.

\section{One structure, four disconnected literatures}

Stand back and the shape is the same in each account: a parametric convex quadratic, an
active set, a piecewise-linear path traced by a ratio test.

Optimization has, all along, had the general language for this---parametric programming,
continuation and homotopy methods, active-set sensitivity analysis---within which every
instance is a special case. In that precise sense none of the device is new. What is
striking is not that the structure existed but that the literatures using it stayed
largely disconnected. The portfolio, regression, kernel, and control versions were
developed, taught, and implemented in parallel, each with its own vocabulary
(Appendix~\ref{app:notation}), its own canonical references, and its own software, and
the general framework that subsumes them was seldom the frame in which any one of them
was presented. A practitioner of one could, and routinely did, work a career without
meeting the others.

It is that disconnection, rather than any depth hidden in the mathematics, that kept the
shared structure from view. And it is a reason to be suspicious, in one's own field, of
any method presented as native to it.

\begin{quote}
That an idea is reached more than once, independently, is less an oversight to correct
than a measure of how good it is.
\end{quote}
